% GENERATED FILE. Edit canonical lessons, then run npm run generate:books.
% canonical-source-hash: fnv1a64:0200d187da79ed9f
% canonical-lessons: SA-C540-asannah, SA-C540-avasistah, SA-C540-samaptih, SA-C540-kalaksepah, SA-C540-ativahayati

\chapter{Time Running Out}
\label{ch:sa-time-running-out}
\hlchaptermodality{540}

\begin{chapteropening}
\textbf{By the end of this chapter:} \emph{I can talk about deadlines and what happens if we are slow.}

Complete the last lesson of chapter 540: I can talk about deadlines and what happens if we are slow.
\end{chapteropening}

\section[\texorpdfstring{\sk{आसन्नः} (āsannaḥ)}{आसन्नः (āsannaḥ)}]{\sk{आसन्नः} (āsannaḥ) — close at hand}

\label{lesson:SA-C540-asannah}



\begin{quote}
Before the new one: say the Sanskrit for suspects, then the Sanskrit for keeps safe.
\end{quote}

\subsection*{You'll want to know: \sk{आसन्नः}}
\textbf{\sk{आसन्नः}} — \emph{āsannaḥ} — \textquotedblleft{}close at hand\textquotedblright{}.

\begin{quote}
\sk{यदि} \sk{परीक्षा} \sk{आसन्ना} \sk{अस्ति}, \sk{तर्हि} \sk{अद्य} \sk{एव} \sk{पठतु।}
\end{quote}

\begin{quote}
\emph{yadi parīkṣā āsannā asti, tarhi adya eva paṭhatu.}
\end{quote}

\begin{quote}
If the exam is close, study today.
\end{quote}

\begin{cousinweb}
\textbf{\sk{आसन्नः}} (\emph{āsannaḥ}) is built on \emph{ā} and the root √sad, \textquotedblleft{}to sit\textquotedblright{}, as its past participle: sat down near; MW gives \emph{āsanna} as seated down, and near, close at hand (MW).
\end{cousinweb}

\begin{grammarlens}[title={What you've built}]
One more word for saying what happens if.
\end{grammarlens}

\subsection*{Your turn}
\begin{itemize}
\raggedright
  \item \emph{Say it:} \emph{āsannaḥ}
  \item \emph{Say it:} \emph{āsannaḥ}, once more
  \item \emph{Say it:} \emph{yadi parīkṣā āsannā asti, tarhi adya eva paṭhatu.} — If the exam is close, study today.
  \item \emph{Recall:} say the Sanskrit for suspects, then the Sanskrit for keeps safe, then say \emph{āsannaḥ} again
\end{itemize}

\begin{tcolorbox}[breakable,colback=teal!4,colframe=teal!35!black,title={Before you move on}]
What does \sk{आसन्नः} mean? (\textquotedblleft{}Close at hand\textquotedblright{}.) Say it once more.
\end{tcolorbox}
\section[\texorpdfstring{\sk{अवशिष्टः} (avaśiṣṭaḥ)}{अवशिष्टः (avaśiṣṭaḥ)}]{\sk{अवशिष्टः} (avaśiṣṭaḥ) — left over}

\label{lesson:SA-C540-avasistah}



\begin{quote}
Before the new one: say the Sanskrit for keeps safe, then the Sanskrit for close at hand.
\end{quote}

\subsection*{You'll want to know: \sk{अवशिष्टः}}
\textbf{\sk{अवशिष्टः}} — \emph{avaśiṣṭaḥ} — \textquotedblleft{}left over\textquotedblright{}.

\begin{quote}
\sk{यदि} \sk{किञ्चित्} \sk{अन्नम्} \sk{अवशिष्टम्} \sk{अस्ति}, \sk{तर्हि} \sk{श्वः} \sk{खादिष्यामः।}
\end{quote}

\begin{quote}
\emph{yadi kiñcit annam avaśiṣṭam asti, tarhi śvaḥ khādiṣyāmaḥ.}
\end{quote}

\begin{quote}
If any food is left over, we'll eat it tomorrow.
\end{quote}

\begin{cousinweb}
\textbf{\sk{अवशिष्टः}} (\emph{avaśiṣṭaḥ}) is built on \emph{ava} and the root √śiṣ, \textquotedblleft{}to leave\textquotedblright{}, as its past participle; MW gives \emph{avaśiṣṭa} as left, remaining (MW).
\end{cousinweb}

\begin{grammarlens}[title={What you've built}]
One more word for saying what happens if.
\end{grammarlens}

\subsection*{Your turn}
\begin{itemize}
\raggedright
  \item \emph{Say it:} \emph{avaśiṣṭaḥ}
  \item \emph{Say it:} \emph{avaśiṣṭaḥ}, once more
  \item \emph{Say it:} \emph{yadi kiñcit annam avaśiṣṭam asti, tarhi śvaḥ khādiṣyāmaḥ.} — If any food is left over, we'll eat it tomorrow.
  \item \emph{Recall:} say the Sanskrit for keeps safe, then the Sanskrit for close at hand, then say \emph{avaśiṣṭaḥ} again
\end{itemize}

\begin{tcolorbox}[breakable,colback=teal!4,colframe=teal!35!black,title={Before you move on}]
What does \sk{अवशिष्टः} mean? (\textquotedblleft{}Left over\textquotedblright{}.) Say it once more.
\end{tcolorbox}
\section[\texorpdfstring{\sk{समाप्तिः} (samāptiḥ)}{समाप्तिः (samāptiḥ)}]{\sk{समाप्तिः} (samāptiḥ) — completion}

\label{lesson:SA-C540-samaptih}



\begin{quote}
Before the new one: say the Sanskrit for close at hand, then the Sanskrit for left over.
\end{quote}

\subsection*{You'll want to know: \sk{समाप्तिः}}
\textbf{\sk{समाप्तिः}} — \emph{samāptiḥ} — \textquotedblleft{}completion\textquotedblright{}.

\begin{quote}
\sk{यदि} \sk{कार्यस्य} \sk{समाप्तिः} \sk{भवति}, \sk{तर्हि} \sk{वयं} \sk{विश्रामं} \sk{करिष्यामः।}
\end{quote}

\begin{quote}
\emph{yadi kāryasya samāptiḥ bhavati, tarhi vayaṃ viśrāmaṃ kariṣyāmaḥ.}
\end{quote}

\begin{quote}
When the work is finished, we will rest.
\end{quote}

A feminine noun.

\begin{cousinweb}
\textbf{\sk{समाप्तिः}} (\emph{samāptiḥ}) is built on \emph{sam}, \textquotedblleft{}wholly\textquotedblright{}, and the root √āp, \textquotedblleft{}to reach\textquotedblright{}, which also gives \emph{samāpayati}, \textquotedblleft{}finishes\textquotedblright{}; MW gives \emph{samāpti} as accomplishment, completion, conclusion (MW).
\end{cousinweb}

\begin{grammarlens}[title={What you've built}]
One more word for saying what happens if.
\end{grammarlens}

\subsection*{Your turn}
\begin{itemize}
\raggedright
  \item \emph{Say it:} \emph{samāptiḥ}
  \item \emph{Say it:} \emph{samāptiḥ}, once more
  \item \emph{Say it:} \emph{yadi kāryasya samāptiḥ bhavati, tarhi vayaṃ viśrāmaṃ kariṣyāmaḥ.} — When the work is finished, we will rest.
  \item \emph{Recall:} say the Sanskrit for close at hand, then the Sanskrit for left over, then say \emph{samāptiḥ} again
\end{itemize}

\begin{tcolorbox}[breakable,colback=teal!4,colframe=teal!35!black,title={Before you move on}]
What does \sk{समाप्तिः} mean? (\textquotedblleft{}Completion\textquotedblright{}.) Say it once more.
\end{tcolorbox}
\section[\texorpdfstring{\sk{कालक्षेपः} (kālakṣepaḥ)}{कालक्षेपः (kālakṣepaḥ)}]{\sk{कालक्षेपः} (kālakṣepaḥ) — a waste of time; delay}

\label{lesson:SA-C540-kalaksepah}



\begin{quote}
Before the new one: say the Sanskrit for left over, then the Sanskrit for completion.
\end{quote}

\subsection*{You'll want to know: \sk{कालक्षेपः}}
\textbf{\sk{कालक्षेपः}} — \emph{kālakṣepaḥ} — \textquotedblleft{}a waste of time; delay\textquotedblright{}.

\begin{quote}
\sk{यदि} \sk{कालक्षेपः} \sk{भवति}, \sk{तर्हि} \sk{वयं} \sk{समये} \sk{न} \sk{आगमिष्यामः।}
\end{quote}

\begin{quote}
\emph{yadi kālakṣepaḥ bhavati, tarhi vayaṃ samaye na āgamiṣyāmaḥ.}
\end{quote}

\begin{quote}
If there is a delay, we will not arrive on time.
\end{quote}

A masculine noun.

\begin{cousinweb}
\textbf{\sk{कालक्षेपः}} (\emph{kālakṣepaḥ}) is built from \emph{kāla}, \textquotedblleft{}time\textquotedblright{}, and \emph{kṣepa}, \textquotedblleft{}a throwing, a letting go\textquotedblright{}, from the root √kṣip, \textquotedblleft{}to throw\textquotedblright{}; MW gives \emph{kālakṣepa} as letting time pass, delay, loss of time (MW).
\end{cousinweb}

\begin{grammarlens}[title={What you've built}]
One more word for saying what happens if.
\end{grammarlens}

\subsection*{Your turn}
\begin{itemize}
\raggedright
  \item \emph{Say it:} \emph{kālakṣepaḥ}
  \item \emph{Say it:} \emph{kālakṣepaḥ}, once more
  \item \emph{Say it:} \emph{yadi kālakṣepaḥ bhavati, tarhi vayaṃ samaye na āgamiṣyāmaḥ.} — If there is a delay, we will not arrive on time.
  \item \emph{Recall:} say the Sanskrit for left over, then the Sanskrit for completion, then say \emph{kālakṣepaḥ} again
\end{itemize}

\begin{tcolorbox}[breakable,colback=teal!4,colframe=teal!35!black,title={Before you move on}]
What does \sk{कालक्षेपः} mean? (\textquotedblleft{}A waste of time; delay\textquotedblright{}.) Say it once more.
\end{tcolorbox}
\section[\texorpdfstring{\sk{अतिवाहयति} (ativāhayati)}{अतिवाहयति (ativāhayati)}]{\sk{अतिवाहयति} (ativāhayati) — spends time}

\label{lesson:SA-C540-ativahayati}



\begin{quote}
Before the new one: say the Sanskrit for completion, then the Sanskrit for a waste of time.
\end{quote}

\subsection*{You'll want to know: \sk{अतिवाहयति}}
\textbf{\sk{अतिवाहयति}} — \emph{ativāhayati} — \textquotedblleft{}spends time\textquotedblright{}.

\begin{quote}
\sk{यदि} \sk{वर्षा} \sk{भवति}, \sk{तर्हि} \sk{वयं} \sk{दिनं} \sk{गृहे} \sk{अतिवाहयामः।}
\end{quote}

\begin{quote}
\emph{yadi varṣā bhavati, tarhi vayaṃ dinaṃ gṛhe ativāhayāmaḥ.}
\end{quote}

\begin{quote}
If it rains, we spend the day at home.
\end{quote}

Past (he …): \sk{सः} \sk{दिनं} \sk{गृहे} \sk{अतिवाहितवान्}, \emph{saḥ dinaṃ gṛhe ativāhitavān}.

\begin{cousinweb}
\textbf{\sk{अतिवाहयति}} (\emph{ativāhayati}) is built on \emph{ati}, \textquotedblleft{}beyond\textquotedblright{}, and the causative of the root √vah, \textquotedblleft{}to carry\textquotedblright{}, which also gives \emph{vahati}; MW gives \emph{ativāhayati} as to let pass, to let time pass, to spend it (MW).
\end{cousinweb}

\begin{grammarlens}[title={What you've built}]
One more word for saying what happens if.
\end{grammarlens}

\subsection*{Your turn}
\begin{itemize}
\raggedright
  \item \emph{Say it:} \emph{ativāhayati}
  \item \emph{Say it:} \emph{ativāhayati}, once more
  \item \emph{Say it:} \emph{yadi varṣā bhavati, tarhi vayaṃ dinaṃ gṛhe ativāhayāmaḥ.} — If it rains, we spend the day at home.
  \item \emph{Recall:} say the Sanskrit for completion, then the Sanskrit for a waste of time, then say \emph{ativāhayati} again
\end{itemize}

\begin{tcolorbox}[breakable,colback=teal!4,colframe=teal!35!black,title={Before you move on}]
What does \sk{अतिवाहयति} mean? (\textquotedblleft{}Spends time\textquotedblright{}.) Say it once more.
\end{tcolorbox}
